% Outline
%
% In this section, we briefly introduce the background of flow-basec/diffusion generative models and recent advances in one-step generation models.
% My plan for this section consists of four paragraphs:
%
% 1. Kick start the introduction with a paragraph discussing about flow-matching and diffusion models. It ends with the limitation of numerical integration.
% 2. Following tha paragraph, describe the recent advances in one-step generation models, including early works on distillation-based one-step generation followed by the introduction of Mean Flows. The paragraph should end by highlighting the limitations: non-decreasing and high-variance loss.
% 3. Following up on our discussion about the limitations, we introduce the primary theoretical contribution of this work: curvature-variance principle, and highlight how we position our contribution amongst concurrent works.
% 4. Finally, we introduce VaMF, a concrete yet simple method to address the analyzed issue and improve the overall performance and training efficiency of the original Mean Flows.

\section{Introduction}
\label{sec: introduction}

%%%%%%%%%%%%%%%%%%%%%%%%%%%%%%%%%%%%%%%%%
% introduction of flow-based/diffusion
Deep generative models that leverage deep neural networks to learn and generate samples from unknown data distributions have achieved profound success. One of the predominant paradigms adopts flow-based generative models~\cite{chen2018neural,rezende2015variational,dinh2017density,grathwohl2019ffjord,lipman2023flowmatchinggenerativemodeling,tong2024improvinggeneralizingflowbasedgenerative,albergo2023buildingnormalizingflowsstochastic} that learn continuous-time transformations between a simple prior and a target data distribution. The existing conditional flow matching~\cite{lipman2023flowmatchinggenerativemodeling,tong2024improvinggeneralizingflowbasedgenerative} and stochastic interpolant~\cite{albergo2023buildingnormalizingflowsstochastic,albergo2025stochasticinterpolants} enable simulation-free training by regressing a velocity field, achieving sample quality competitive with diffusion models~\cite{sohldickstein2015deep,song2019generative,ho2020denoising,nichol2021improved,song2021scorebased,dhariwal2021diffusion,karras2022elucidating} and variational autoencoders~\cite{kingma2014auto}. However, flow-based models share the same limitation as diffusion models as they require iterative integration at inference time, which can suffer from extensive computational costs and numerical instability.

%%%%%%%%%%%%%%%%%%%%%%%%%%%%%%%%%%%%%%%%%
% few-/one-step generation models
The issue motivates consecutive studies on reducing the number of function evaluations to one or few steps~\cite{song2021denoising}. Existing approaches include progressive distillation~\cite{salimans2022progressivedistillationfastsampling}, distribution matching~\cite{yin2024distributionmatchingdistillation,yin2024improveddistributionmatchingdistillation}, consistency models~\cite{song2023consistencymodels,lu2025simplifyingstabilizingscalingcontinuoustime}, and flow map matching~\cite{boffi2025flowmapmatchingstochastic}, but most require a pre-trained teacher model. Recently, Mean Flows~\cite{geng2025meanflowsonestepgenerative} offers an appealing \textit{distillation-free} alternative by learning a two-parameter average velocity field $\vu(\vx,r,t)$ through a total-derivative identity between the average and marginal velocity field. In principle, this identity provides exact self-supervision. The compound prediction $\vu+(t-r)\frac{\mathrm{d}}{\mathrm{d}t}\vu$ should match the instantaneous velocity at convergence, requiring no external teacher. Nevertheless, empirical training of Mean Flows is plagued by non-decreasing loss and prohibitively high gradient variance~\cite{zhang2025alphaflowunderstandingimprovingmeanflow,geng2025improvedmeanflowschallenges}.

%%%%%%%%%%%%%%%%%%%%%%%%%%%%%%%%%%%%%%%%%
% our contributions
In this work, we conduct a theoretical analysis that attributes the issue to replacing the marginal velocity fields by conditional fields in the total derivatives and the regression target, and establish the \textbf{curvature-variance principle}. Specifically, we prove that: \textbf{(i)} the per-step gradient variance is amplified by a Jacobi factor $\Tr(\mJ\Sigma_{\vv'}\mJ^{\!\top})$ where $\mJ=(t-r)\partial_{\vz}\vu_{\vtheta}-\mI_{d}$, and the stop-gradient operator creates a semi-gradient gap that blinds the optimizer to this amplification; \textbf{(ii)} there exists a connection between the Jacobi and the curvature gap $\Delta:=\vu-\vv$, which satisfies $\Delta\approx-\frac{(t-r)}{2}\frac{\mathrm{D}\vv}{\mathrm{D}t}$, so it is the local flow acceleration that causes the average and instantaneous velocities to diverge; and \textbf{(iii)} the acceleration-to-noise ratio $\text{ANR}=(t-r)\|\frac{\mathrm{D}\vv}{\mathrm{D}t}\|/\sqrt{\Tr(\Sigma_{\vv'})}$ governs the optimal sampling distribution over intervals, connecting scheduling heuristics in prior work to a principled criterion.

Building on this analysis, we propose \textbf{Variance-Aware MeanFlow (VaMF)}, a backbone-preserving training framework that addresses each diagnostic in turn. First, an \textit{EMA velocity anchor} replaces the stochastic JVP tangent with a deterministic estimate, provably eliminating the Jacobi-amplified gradient noise. Then, a \textit{flow-matching anchor} supervises the boundary condition at $r\!\to\!t$, removing the transient bias that an inexact tangent would otherwise introduce. Finally, a closed-form \textit{trace weight} $w\!\propto\!1/(1+\sigma_{t}^{2}\Tr(\mJ\mJ^{\!\top})/d)$ downweights each sample by the local Jacobian-deviation norm. VaMF preserves the original Mean Flow MSE objective and adds only one Jacobian-vector product per step. Our contributions are as follows:
\begin{description}
    \item[\namedlabel{cb: thm}{C1}] We establish the curvature-variance principle, which attributes the non-decreasing, high-variance loss of Mean Flows to a single mechanism---the Jacobi-amplified per-step gradient variance and its connection to the local flow acceleration---and provide the first formal quantification of the semi-gradient gap induced by the stop-gradient operator.
    \item[\namedlabel{cb: vamf}{C2}] We derive VaMF, a backbone-preserving training framework that combines an EMA velocity anchor, a flow-matching anchor, and a Hutchinson-estimated trace weight to eliminate variance amplification and the associated tangent bias.
    \item[\namedlabel{cb: exp}{C3}] Empirically, VaMF improves Sliced Wasserstein-$1$ ($\text{SW}_{1}$) over vanilla MeanFlow by $-23\%$, $-13\%$, $-23\%$, and $-54\%$ on the four 2D datasets, scales favorably from $d{=}2$ up to $d{=}64$ on the dense-Gaussian-mixture sweep, and transfers to large training of latent-space image generation models on the ImageNet datasets.
\end{description}
